\begin{abstract}
Recent generative image models such as Stable Diffusion allow a publicly released photograph, portraits in particular, to be modified with a text prompt, which raises privacy and security concerns. Image immunization addresses this threat by modifying a photograph before its release. Most existing methods superimpose an additive adversarial perturbation on the image, which reduces the chance that an edit succeeds. The energy of this perturbation lies mainly at high spatial frequencies. Purification operations such as JPEG compression, Gaussian blurring, and encoder-based reconstruction therefore remove it with little effort. This vulnerability stems from the spatial structure of the perturbation. We propose color-mapping immunization, a filter-like defense that encodes the protection in the colors of the image. Each image receives its own 3D color lookup table (LUT), the operator of global color grading. Acting on colors alone, the LUT introduces no spatial structure into the image, and the protection it carries persists through purification. The protected image appears as a color-graded version of the original. On InstructPix2Pix at matched distortion, the edit displacement retained after purification, averaged over all purification methods, is \NUM{} times that of additive defenses.
\end{abstract}
